\scanpage{152}\mspage{7}%
A) Le formalisme des $Rf_!, f^!$ (combiné à $Rf_*, Lf^{\unsure{\alpha|*}}$, $\overset{\mathbb{L}}{\otimes}$ et $\mathbb{R}\mathrm{Hom}$, ``les six opérations'') porte implicitement en son sein la définition de l'homologie et l'identité essentielle entre \unsure{fonctions} homologiques et cohomologiques. \deleted{\unsure{E}} On dispose pour le moment de ce formalisme pour les faisceaux quasi-cohérents sur les schémas -- séminaire Hartshorne -- pour les espaces top. et les faisceaux de coeff. arbitraires dessus -- Verdier, exposé Bourbaki -- et en coh. étale des schémas pour les coefficients ``discrets'' \inserted{(de torsion ou ``$l$-adiques'')} \uline{premiers aux \deleted{\unsure{la|le}} car. résiduelles} -- SGA 5 -- \inserted{enfin pour les faisc. coh. sur les espaces analytiques -- Verdier-Ruget --}\,; (Le formalisme reste à développer dans le contexte cristallin, et en car. $0$ dans le contexte des modules stratifiés \inserted{à singularités} à la Deligne, avec éventuellement -- sur le corps $\mathbb{C}$ -- des structures de Hodge à la clef, enfin dans le contexte des motifs ; je suis persuadé qu'il est présent un peu partout encore -- partout, peut-être, où il y a un formalisme de nature cohomologique.) Travaillant en coh. étale sur un schéma \inserted{\unsure{séparé}} de t.f. sur un corps $k$, disons, avec un anneau de coeff. $\Lambda$ de torsion premier à la \unsure{car.|car}\unsure{,}
% archive sheet number: 152
